\section{Research Methodology}

\subsection{Research design}
% TODO

\subsection{Tool selection strategy}
The selection of risk managment tools were mostly based on the requirements detailed in the assignment description. Most of the tools used in this project were the ones already identified in the assignment description, with additional tools found through independent research. 

The independent research into other tools not found in the assignment description consisted of using search engines such as google.com in order to find tools relevant for this project. The searches included terms such as "risk managment tool" and "risk analysis tools". After finding a tool, the research of said tool begins. This consists of finding reviews for it through review websites and other sites such as reddit.com to get a broad base of opinions.

Only tools relating to risk analysis was chosen for this Project. The selection was purposeful. Not every tool could be chosen since there are a vast amount available online, but availability was prioritized. 


\subsection{Development of Evaluation Criteria}
The development of evaluation criteria used in this project was developed through a mixture of different risk managment frameworks such as ISO/IEC 27005 and OCTAVE. The reason it was done this way was to get a broad criteria that could be used in evaluating the tools in different hypothetical cases. Combining these frameworks gives us an overview of which comparisons to make regarding the risk analysis tools. These comparisons include core activities such as risk identification, analysis method, documentation, how capable the tool is in reporting incidents, how broad and detailed the risk mitigations are and how closely it aligns with established risk assessment protocols.


\subsection{Evaluation Approach}
The evaluation approach consists of testing the tools in different hypothethical situations regarding small to large businesses. These scenarios were generated using large Language models such as ChatGPT. LLMs were used in order to speed up the process of creating cases for businesses differing in size. The tools themselves and how they perform compared to established protocols was the focus of this project, so speeding up the case-building for those tools seemed natural.

The reason for choosing businesses varying in size for these hypothethical situations to test the tools in, was to simulate how these tools perform under different organizational complexity, difference in the availabel resources, difference in risk tolerance, risk exposure and difference in technical understanding regarding risk analysis.

The tools were assessed on how well they supported core risk managment protocols and processes and how well they followed established frameworks. The ease of use was also in focus. Other criteria such as scalability and adaptability to different organizational sizes were also documented.

\subsection{Limitations and Validity Considerations}
This project has several limitations. The first limitation is that the evaluation is a limited hands-on testing rather than a long term implementation that can be found in real organizations. The testing happens under predetermined scenarios, and while we have tried our best to implement multiple real-world scenarios, it will be different to how they may perform in an actual organization which tends to be much more dynamic.

The second limitation is that the evaluation is conducted using hypothetical organizational scenarios. The scenarios have been made to be as realistic as possible, but organizatiosn are much more diverse than those scenarios, things like work culture, organizational structure, and current risk landscapes may vary. Because of this, the findings in this project should not be thought of as definitive. 

The third limitation is that only a certain amount of tools were chosen for this project. There is a vast amount of tools to choose from online, and some of those tools might be superior or inferior in different aspects compared to the ones chosen for this project. The risk management landscape might also change over time, which in turn might make the tools more effective and with higher capabilities in the future.

Regarding validity, the criteria to evaluate the tools are based on established frameworks withing the risk management landscape. An objective approach has been taken in this project, but there are also some unavoidable subjective assessments when it comes to judging how complicated the tools are to use.
